\ifdefined\inMainDoc\else
\documentclass[12pt,a4paper]{article}
% ================================================================
%  PRÉAMBULE COMMUN - ANNALES CORRIGÉES
%  Université de Kara - FaST - Licence Fondamentale MPC
%  Design optimisé impression noir et blanc
% ================================================================

% -- ENCODAGE & LANGUE --
\usepackage[utf8]{inputenc}
\usepackage[T1]{fontenc}
\usepackage[french]{babel}
\usepackage{lmodern}
\usepackage{microtype}

% -- MISE EN PAGE --
\usepackage[a4paper, top=2.8cm, bottom=2.5cm, left=2.5cm, right=2.5cm, headheight=14pt]{geometry}
\usepackage{setspace}
\setstretch{1.2}
\usepackage{parskip}
\setlength{\parindent}{1.2em}
\setlength{\parskip}{4pt}

% -- MATHS --
\usepackage{amsmath, amssymb, mathtools, amsthm}
\usepackage{mathrsfs}
\usepackage{dsfont}
\usepackage{bm}
\usepackage{cancel}
\usepackage{textcomp}
% Définition de \oiint (intégrale de surface fermée) sans package supplémentaire
\newcommand{\oiint}{\mathop{\iint\!\!\!\!\!\!\!\bigcirc}\limits}

% -- PHYSIQUE & CHIMIE --
\usepackage{physics}
\usepackage{siunitx}
% Lorsque physics est chargé, siunitx omet \qty. On le restaure ici :
\AtBeginDocument{\RenewCommandCopy\qty\SI}
\sisetup{locale = FR, detect-all, output-decimal-marker={,}}
% Commande de substitution pour les unités posant problème avec pdflatex
% (ohm, degreeCelsius, micro, angstrom, percent utilisent des caractères non-ASCII)
\newcommand{\qtyohm}[1]{\ensuremath{\num{#1}\,\Omega}}
\newcommand{\qtydegc}[1]{\ensuremath{\num{#1}\,{}^\circ\text{C}}}
\newcommand{\qtyangstrom}[1]{\ensuremath{\num{#1}\,\text{\AA}}}
\newcommand{\qtypercent}[1]{\ensuremath{\num{#1}\,\%}}
\newcommand{\qtyum}[1]{\ensuremath{\num{#1}\,\mu\text{m}}}
\usepackage[version=4]{mhchem}

% -- GRAPHISMES --
\usepackage{graphicx}
\usepackage{float}
\usepackage{caption}
\usepackage{subcaption}
\usepackage{wrapfig}
\usepackage{tikz}
\usetikzlibrary{calc, arrows.meta, patterns, shapes.geometric}
\usepackage{pgfplots}
\pgfplotsset{compat=1.18}

% -- TABLEAUX --
\usepackage{booktabs}
\usepackage{array}
\usepackage{tabularx}
\usepackage{multirow}

% -- LISTES --
\usepackage{enumitem}
\setlist[enumerate]{topsep=3pt, itemsep=2pt, parsep=0pt}
\setlist[itemize]{topsep=3pt, itemsep=2pt, parsep=0pt, label={.}}

% -- BOÎTES (noir et blanc) --
\usepackage{mdframed}
\usepackage{tcolorbox}
\tcbuselibrary{skins, breakable, theorems}

% Boîte EXERCICE : encadré avec filet épais, fond blanc
\newmdenv[
  linewidth=1.2pt,
  linecolor=black,
  roundcorner=3pt,
  skipabove=10pt,
  skipbelow=6pt,
  innertopmargin=8pt,
  innerbottommargin=8pt,
  innerleftmargin=10pt,
  innerrightmargin=10pt,
  backgroundcolor=white,
  frametitle={\bfseries\small Exercice},
  frametitlealignment=\raggedright,
  frametitleaboveskip=4pt,
  frametitlebelowskip=4pt,
]{boiteexercice}

% Boîte CORRIGÉ : fond gris très clair + filet fin
\newmdenv[
  linewidth=0.6pt,
  linecolor=black,
  leftline=true,
  rightline=false,
  topline=false,
  bottomline=false,
  linecolor=black,
  innerleftmargin=12pt,
  innerrightmargin=8pt,
  innertopmargin=6pt,
  innerbottommargin=6pt,
  backgroundcolor=gray!8,
  skipabove=4pt,
  skipbelow=8pt,
]{boitecorrige}

% Boîte RÉSUMÉ DE COURS : double encadré
\newmdenv[
  linewidth=0.8pt,
  linecolor=black,
  roundcorner=2pt,
  backgroundcolor=gray!12,
  innertopmargin=10pt,
  innerbottommargin=10pt,
  innerleftmargin=12pt,
  innerrightmargin=12pt,
  skipabove=12pt,
  skipbelow=8pt,
]{boiteresume}

% Boîte ATTENTION/REMARQUE : barre gauche épaisse
\newmdenv[
  leftline=true,
  rightline=false,
  topline=false,
  bottomline=false,
  linewidth=3pt,
  linecolor=black,
  backgroundcolor=gray!5,
  innerleftmargin=12pt,
  innerrightmargin=8pt,
  innertopmargin=5pt,
  innerbottommargin=5pt,
  skipabove=6pt,
  skipbelow=6pt,
]{boiteremarque}

% -- DIFFICULTÉ DES EXERCICES --
\newcommand{\facile}{$\star$}
\newcommand{\moyen}{$\star\!\star$}
\newcommand{\difficile}{$\star\!\star\!\star$}
\newcommand{\niveau}[1]{\hfill\textbf{#1}\hspace{2pt}}

% -- EN-TÊTES ET PIEDS DE PAGE --
\usepackage{fancyhdr}
\pagestyle{fancy}
\fancyhf{}
\renewcommand{\headrulewidth}{0.5pt}
\renewcommand{\footrulewidth}{0.3pt}

% -- HYPERLIENS --
\usepackage[hidelinks]{hyperref}
\hypersetup{
    colorlinks=false,
    pdftitle={Annale Corrigée - Licence MPC - Université de Kara}
}

% -- TITRES --
\usepackage{titlesec}
\titleformat{\section}{\Large\bfseries}{\thesection.}{0.8em}{}[\vspace{2pt}\hrule\vspace{4pt}]
\titleformat{\subsection}{\large\bfseries}{\thesubsection.}{0.7em}{}
\titleformat{\subsubsection}{\normalsize\bfseries}{\thesubsubsection.}{0.6em}{}

% -- THÉORÈMES --
\theoremstyle{definition}
\newtheorem{defn}{Définition}[section]
\newtheorem{prop}{Propriété}[section]
\newtheorem{thm}{Théorème}[section]
\newtheorem{corol}{Corollaire}[section]
\newtheorem{exemple}{Exemple}[section]

% -- COMMANDES MATHÉMATIQUES UTILES --
\newcommand{\vect}[1]{\overrightarrow{#1}}
\newcommand{\R}{\mathbb{R}}
\newcommand{\N}{\mathbb{N}}
\newcommand{\Z}{\mathbb{Z}}
\newcommand{\C}{\mathbb{C}}
\renewcommand{\d}{\mathrm{d}}
\newcommand{\deriv}[2]{\dfrac{\d #1}{\d #2}}
\newcommand{\dpart}[2]{\dfrac{\partial #1}{\partial #2}}
\newcommand{\rot}{\overrightarrow{\mathrm{rot}}\,}
\renewcommand{\div}{\mathrm{div}\,}
\renewcommand{\grad}{\overrightarrow{\mathrm{grad}}\,}
\newcommand{\e}{\mathrm{e}}
\newcommand{\ii}{\mathrm{i}}
\renewcommand{\Tr}{\mathrm{Tr}}

% Numérotation des équations par section
\numberwithin{equation}{section}

% -- RÉSULTATS ENCADRÉS --
\newcommand{\res}[1]{\[\boxed{#1}\]}

% -- REMARQUE INLINE --
\newcommand{\rmq}[1]{\begin{boiteremarque}\textbf{Remarque.} #1\end{boiteremarque}}

% -- MISE EN VALEUR DU RÉSULTAT FINAL --
\newcommand{\reponse}[1]{%
  \begin{center}
    \fbox{\begin{minipage}{0.92\textwidth}\centering #1\end{minipage}}
  \end{center}%
}


\lhead{\small\textit{Électromagnétisme dans la Matière - Licence 3}}
\rhead{\small\textit{FaST, Université de Kara}}
\cfoot{\thepage}

\begin{document}
\fi

\begin{center}
  {\LARGE\bfseries ÉLECTROMAGNÉTISME DANS LA MATIÈRE}\\[0.3cm]
  {\normalsize Licence 3 - Physique}\\[0.1cm]
  \rule{0.7\textwidth}{0.8pt}
\end{center}

\section*{Résumé du cours}

\begin{boiteresume}

\textbf{1. Polarisation diélectrique.}
Dans un diélectrique LHI (linéaire, homogène, isotrope), les dipôles microscopiques créent une polarisation macroscopique $\vec{P} = \epsilon_0\chi_e\vec{E}$, avec $\chi_e = \epsilon_r-1$ la susceptibilité électrique. Le vecteur déplacement électrique est :
\[
\vec{D} = \epsilon_0\vec{E} + \vec{P} = \epsilon_0\epsilon_r\vec{E}.
\]
Les charges de polarisation volumiques et surfaciques sont $\rho_p = -\div\vec{P}$ et $\sigma_p = \vec{P}\cdot\hat{n}$.

\textbf{2. Champ local (relation de Clausius-Mossotti).}
Pour un diélectrique, le champ local ressenti par un dipôle est :
\[
\vec{E}_\text{loc} = \vec{E} + \frac{\vec{P}}{3\epsilon_0},
\]
ce qui conduit, pour une polarisabilité atomique $\alpha$, à la relation de Clausius-Mossotti :
$\dfrac{\epsilon_r-1}{\epsilon_r+2} = \dfrac{n\alpha}{3\epsilon_0}$, où $n$ est la densité de dipôles.

\textbf{3. Magnétisation.}
Dans un milieu magnétique LHI, $\vec{M} = \chi_m\vec{H}$ avec $\chi_m$ la susceptibilité magnétique. Le champ $\vec{B}$ et le vecteur excitation $\vec{H}$ sont liés par :
\[
\vec{B} = \mu_0(\vec{H}+\vec{M}) = \mu_0\mu_r\vec{H}.
\]
Les courants d'aimantation volumiques et surfaciques sont $\vec{j}_m = \rot\vec{M}$ et $\vec{K}_m = \vec{M}\times\hat{n}$.

\textbf{4. Équations de Maxwell dans la matière.}
\[
\div\vec{D} = \rho_\ell, \quad
\div\vec{B} = 0, \quad
\rot\vec{E} = -\frac{\partial\vec{B}}{\partial t}, \quad
\rot\vec{H} = \vec{j}_\ell + \frac{\partial\vec{D}}{\partial t}.
\]

\textbf{5. Conditions aux limites à une interface.}
Pour une interface entre deux milieux (1) et (2), avec la normale $\hat{n}$ de (1) vers (2) :
\[
(D_2-D_1)\cdot\hat{n} = \sigma_\ell, \quad
(B_2-B_1)\cdot\hat{n} = 0, \quad
(E_2-E_1)\times\hat{n} = \vec{0}, \quad
(H_2-H_1)\times\hat{n} = \vec{K}_\ell.
\]
En l'absence de charges et courants libres surfaciques : la composante normale de $\vec{D}$ et la composante tangentielle de $\vec{E}$ sont continues.

\textbf{6. Ondes électromagnétiques dans les milieux matériels.}
Dans un milieu LHI non absorbant d'indices $\epsilon_r$ et $\mu_r$, l'onde plane se propage avec la vitesse $v = c/n$, où l'indice de réfraction vaut $n = \sqrt{\epsilon_r\mu_r}$. En présence d'absorption, $n$ devient complexe : $\tilde{n} = n + \ii\kappa$ (avec $\kappa$ le coefficient d'extinction). L'intensité décroît en $\e^{-2\kappa\omega z/c}$ (loi de Beer-Lambert). À une interface, les lois de Snell-Descartes s'appliquent : $n_1\sin\theta_1 = n_2\sin\theta_2$.

\end{boiteresume}

\newpage
\subsection*{Exercice 1 - Champ local dans un diélectrique \niveau{\moyen}}

\begin{boiteexercice}
On considère un diélectrique LHI de polarisation uniforme $\vec{P}$ dans un champ extérieur $\vec{E}_0$. On creuse une cavité sphérique de rayon $R$ (très petit à l'échelle macroscopique) centrée sur un point $M$ du diélectrique.
\begin{enumerate}
  \item Exprimer le champ local $\vec{E}_\text{loc}$ en $M$ en fonction du champ macroscopique $\vec{E}$ et du champ $\vec{E}_p'(M)$ créé par les charges de polarisation sur la surface $S'$ de la cavité.
  \item Calculer $\vec{E}_p'(M)$ en fonction de $\vec{P}$ et $\epsilon_0$. En déduire $\vec{E}_\text{loc}$ en fonction de $\vec{P}$, $\epsilon_0$ et $\vec{E}$.
\end{enumerate}
\end{boiteexercice}

\begin{boitecorrige}

\textbf{1. Expression du champ local.}

Le champ en $M$ résulte de trois contributions :
\[
\vec{E}_\text{loc} = \vec{E}_0 + \vec{E}_p + \vec{E}_p',
\]
où $\vec{E}_p$ est le champ créé par les charges de polarisation sur la surface extérieure $S$ du diélectrique, et $\vec{E}_p'$ est celui créé par les charges sur la surface $S'$ de la cavité. Or $\vec{E}_0 + \vec{E}_p = \vec{E}$ (champ macroscopique), d'où :
\[
\vec{E}_\text{loc} = \vec{E} + \vec{E}_p'.
\]

\textbf{2. Calcul de $\vec{E}_p'$.}

La densité surfacique de charges de polarisation sur $S'$ est $\sigma_p = \vec{P}\cdot(-\hat{e}_r) = -P\cos\theta$ (la normale extérieure à la cavité est $-\hat{e}_r$). Le champ élémentaire en $M$ créé par un élément de surface $\d S' = 2\pi R^2\sin\theta\,\d\theta$ est, par la loi de Coulomb :
\[
\d\vec{E}_p' = \frac{1}{4\pi\epsilon_0}\,\frac{\sigma_p\,\d S'}{R^2}\,(-\hat{e}_r).
\]
Par symétrie de révolution autour de l'axe $\vec{P}$, seule la composante sur cet axe subsiste :
\[
E_{p,z}' = \frac{P}{2\epsilon_0}\int_0^{\pi}\cos^2\theta\sin\theta\,\d\theta
= \frac{P}{2\epsilon_0}\cdot\frac{2}{3} = \frac{P}{3\epsilon_0}.
\]
Donc :
\[
\vec{E}_p' = \frac{\vec{P}}{3\epsilon_0},
\qquad\text{et}\qquad
\boxed{\vec{E}_\text{loc} = \vec{E} + \frac{\vec{P}}{3\epsilon_0}.}
\]
Ce résultat est la \emph{relation de Lorentz} pour le champ local.
\end{boitecorrige}

\subsection*{Exercice 2 - Condensateur plan avec deux diélectriques en parallèle \niveau{\moyen}}

\begin{boiteexercice}
Un condensateur plan est formé de deux plaques rectangulaires de dimensions $(a, e)$ séparées d'une distance $d$. Il porte une distribution surfacique de charges libres $\sigma_\ell$. On néglige les effets de bord.
\begin{enumerate}
  \item Condensateur à air. Calculer le champ électrique entre les armatures et la capacité $C_0$.
  \item Le condensateur est maintenant rempli de deux diélectriques de permittivités relatives $\epsilon_{r1}$ et $\epsilon_{r2}$, disposés en parallèle sur des longueurs $e_1$ et $e_2$ respectivement ($e_1+e_2=e$).
  \begin{enumerate}[label=\alph*)]
    \item Calculer les champs $\vec{D}_i$ et $\vec{E}_i$ dans chaque diélectrique $i$.
    \item En déduire une relation entre la densité de charges libres $\sigma_\ell$ et la densité de charges de polarisation $\sigma_p^i$ dans le diélectrique $i$.
    \item Calculer la nouvelle capacité $C$ en fonction de $a$, $e_1$, $e_2$, $d$, $\epsilon_{r1}$ et $\epsilon_{r2}$. Conclure.
    \item Que vaut $C$ si les deux diélectriques sont identiques ?
  \end{enumerate}
\end{enumerate}
\end{boiteexercice}

\begin{boitecorrige}

\textbf{1. Condensateur à air.}

Par le théorème de Gauss, chaque armature crée $\sigma_\ell/(2\epsilon_0)$. Les deux contributions s'ajoutent entre les armatures :
\[
\vec{E}_0 = \frac{\sigma_\ell}{\epsilon_0}\,\hat{e}_z.
\]
La tension est $U_0 = E_0\,d$, et la capacité :
\[
C_0 = \frac{Q_\ell}{U_0} = \frac{\sigma_\ell\,ae}{E_0\,d} = \frac{\epsilon_0\,ae}{d}.
\]

\textbf{2. Avec deux diélectriques en parallèle.}

\textbf{a)} Les diélectriques étant disposés parallèlement aux armatures (même tension à leurs bornes), la continuité de la composante normale de $\vec{D}$ à l'interface armature/diélectrique donne :
\[
\vec{D}_i = \sigma_\ell\,\hat{e}_z
\qquad\Rightarrow\qquad
\vec{E}_i = \frac{\vec{D}_i}{\epsilon_0\epsilon_{ri}} = \frac{\sigma_\ell}{\epsilon_0\epsilon_{ri}}\,\hat{e}_z.
\]

\textbf{b)} La polarisation dans le diélectrique $i$ est $\vec{P}_i = (\epsilon_{ri}-1)\epsilon_0\vec{E}_i = \left(1-\dfrac{1}{\epsilon_{ri}}\right)\sigma_\ell\,\hat{e}_z$.

La charge de polarisation surfacique à l'interface avec l'armature est $\sigma_p^i = P_i\cdot(-1)$ (normale sortante opposée) :
\[
\sigma_p^i = -\left(1 - \frac{1}{\epsilon_{ri}}\right)\sigma_\ell = \left(\frac{1}{\epsilon_{ri}}-1\right)\sigma_\ell.
\]

\textbf{c)} Le système est équivalent à deux condensateurs en parallèle, chacun de surface $ae_i$ et de distance $d$. La capacité totale est :
\[
C = C_1 + C_2 = \frac{\epsilon_0\epsilon_{r1}\,ae_1}{d} + \frac{\epsilon_0\epsilon_{r2}\,ae_2}{d}
= \frac{\epsilon_0\,a}{d}\!\left(\epsilon_{r1}\,e_1+\epsilon_{r2}\,e_2\right).
\]
Deux condensateurs en parallèle ont la même tension à leurs bornes et leurs capacités s'additionnent.

\textbf{d)} Si $\epsilon_{r1}=\epsilon_{r2}=\epsilon_r$, alors $e_1+e_2=e$ et :
\[
C = \frac{\epsilon_0\epsilon_r\,ae}{d} = \epsilon_r C_0.
\]
On retrouve bien que remplir un condensateur d'un diélectrique uniforme multiplie sa capacité par $\epsilon_r$.

\end{boitecorrige}

\subsection*{Exercice 3 - Réfraction à l'interface entre deux diélectriques \niveau{\moyen}}

\begin{boiteexercice}
Une onde électromagnétique plane se propage dans le vide et atteint une interface plane avec un diélectrique de permittivité relative $\epsilon_r = 4$ (milieu non magnétique). L'angle d'incidence est $\theta_i = 30°$.
\begin{enumerate}
  \item Calculer l'indice de réfraction $n$ du diélectrique.
  \item Appliquer la loi de Snell-Descartes pour trouver l'angle de réfraction $\theta_r$.
  \item Existe-t-il un angle de Brewster pour une polarisation TM ? Si oui, le calculer.
\end{enumerate}
\end{boiteexercice}

\begin{boitecorrige}

\textbf{1.} Pour un milieu non magnétique ($\mu_r=1$) : $n = \sqrt{\epsilon_r} = \sqrt{4} = 2$.

\textbf{2.} La loi de Snell-Descartes donne ($n_1=1$ pour le vide, $n_2=2$) :
\[
n_1\sin\theta_i = n_2\sin\theta_r \implies \sin\theta_r = \frac{\sin 30°}{2} = \frac{1}{4}
\implies \theta_r = \arcsin\!\left(\frac{1}{4}\right) \approx 14{,}5°.
\]

\textbf{3.} L'angle de Brewster $\theta_B$ pour la polarisation TM vérifie $\tan\theta_B = n_2/n_1 = 2$, soit :
\[
\theta_B = \arctan(2) \approx 63{,}4°.
\]
À cet angle d'incidence, la composante TM du champ réfléchi s'annule.
\end{boitecorrige}

\subsection*{Exercice 4 - Conditions aux limites magnétiques \niveau{\moyen}}

\begin{boiteexercice}
On considère une interface plane séparant deux milieux magnétiques LHI de perméabilités relatives $\mu_{r1}$ et $\mu_{r2}$. Un champ $\vec{B}_1$ fait un angle $\alpha_1$ avec la normale à l'interface dans le milieu 1.
\begin{enumerate}
  \item Écrire les conditions aux limites sur $\vec{B}$ et $\vec{H}$ à l'interface (en l'absence de courant surfacique libre).
  \item Établir la loi de réfraction du champ magnétique : $\dfrac{\tan\alpha_1}{\tan\alpha_2} = \dfrac{\mu_{r1}}{\mu_{r2}}$.
\end{enumerate}
\end{boiteexercice}

\begin{boitecorrige}

\textbf{1. Conditions aux limites.}

En l'absence de courant surfacique libre, les conditions aux limites sont :
\begin{align*}
B_{2n} &= B_{1n} \quad\text{(composante normale de $\vec{B}$, continue)}\\
H_{2t} &= H_{1t} \quad\text{(composante tangentielle de $\vec{H}$, continue)}
\end{align*}

\textbf{2. Loi de réfraction.}

Les composantes sont :
\[
B_{1n} = B_1\cos\alpha_1,\quad B_{1t} = B_1\sin\alpha_1,\quad H_{1t} = \frac{B_{1t}}{\mu_0\mu_{r1}} = \frac{B_1\sin\alpha_1}{\mu_0\mu_{r1}}.
\]
La continuité de $B_n$ donne $B_1\cos\alpha_1 = B_2\cos\alpha_2$.

La continuité de $H_t$ donne $\dfrac{B_1\sin\alpha_1}{\mu_{r1}} = \dfrac{B_2\sin\alpha_2}{\mu_{r2}}$.

En divisant la seconde équation par la première :
\[
\frac{\tan\alpha_1}{\mu_{r1}} = \frac{\tan\alpha_2}{\mu_{r2}}
\implies \boxed{\frac{\tan\alpha_1}{\tan\alpha_2} = \frac{\mu_{r1}}{\mu_{r2}}.}
\]
\end{boitecorrige}

\subsection*{Exercice 5 - Onde électromagnétique dans un milieu absorbant \niveau{\difficile}}

\begin{boiteexercice}
On considère une onde électromagnétique plane de pulsation $\omega$ se propageant dans un milieu conducteur de permittivité $\epsilon_r$, de perméabilité $\mu_r = 1$ et de conductivité $\sigma$.
\begin{enumerate}
  \item Montrer que l'équation de propagation pour le champ $\vec{E}$ s'écrit :
  \[
    \Delta\vec{E} = \mu_0\epsilon_0\epsilon_r\frac{\partial^2\vec{E}}{\partial t^2} + \mu_0\sigma\frac{\partial\vec{E}}{\partial t}.
  \]
  \item En cherchant une solution en onde plane $\vec{E} = \vec{E}_0\,\e^{\ii(\tilde{k}z-\omega t)}$, montrer que le vecteur d'onde complexe vérifie $\tilde{k}^2 = \mu_0\epsilon_0\omega^2\!\left(\epsilon_r + \dfrac{\ii\sigma}{\epsilon_0\omega}\right)$.
  \item Identifier la profondeur de peau $\delta$ (distance de pénétration) dans la limite des bons conducteurs ($\sigma\gg\epsilon_0\epsilon_r\omega$).
\end{enumerate}
\end{boiteexercice}

\begin{boitecorrige}

\textbf{1. Équation de propagation.}

La loi d'Ohm donne $\vec{j} = \sigma\vec{E}$. Les équations de Maxwell dans le milieu ($\div\vec{E}=0$ pour une onde transverse) donnent :
\[
\rot(\rot\vec{E}) = -\nabla^2\vec{E} = -\mu_0\frac{\partial}{\partial t}(\rot\vec{H}) = -\mu_0\frac{\partial}{\partial t}\!\left(\sigma\vec{E}+\epsilon_0\epsilon_r\frac{\partial\vec{E}}{\partial t}\right).
\]
D'où l'équation de propagation :
\[
\Delta\vec{E} = \mu_0\epsilon_0\epsilon_r\frac{\partial^2\vec{E}}{\partial t^2} + \mu_0\sigma\frac{\partial\vec{E}}{\partial t}.
\]

\textbf{2. Relation de dispersion.}

En substituant $\vec{E} = \vec{E}_0\,\e^{\ii(\tilde{k}z-\omega t)}$, on obtient $-\tilde{k}^2 = -\mu_0\epsilon_0\epsilon_r\omega^2 + \ii\mu_0\sigma\omega$, soit :
\[
\tilde{k}^2 = \mu_0\epsilon_0\omega^2\!\left(\epsilon_r + \frac{\ii\sigma}{\epsilon_0\omega}\right).
\]
C'est la permittivité complexe $\tilde{\epsilon}_r = \epsilon_r + \ii\sigma/(\epsilon_0\omega)$.

\textbf{3. Profondeur de peau.}

Dans la limite $\sigma\gg\epsilon_0\epsilon_r\omega$ (bon conducteur) :
\[
\tilde{k}^2 \approx \ii\mu_0\sigma\omega
\implies \tilde{k} \approx \frac{1+\ii}{\sqrt{2}}\sqrt{\mu_0\sigma\omega}.
\]
La partie imaginaire donne l'amortissement : l'amplitude décroît comme $\e^{-\Im(\tilde{k})z}$. La profondeur de peau est :
\[
\boxed{\delta = \frac{1}{\Im(\tilde{k})} = \sqrt{\frac{2}{\mu_0\sigma\omega}}.}
\]
L'onde est atténuée d'un facteur $1/e$ sur une distance $\delta$. Pour le cuivre à $50\;\text{Hz}$, $\delta\approx 9\;\text{mm}$.

\end{boitecorrige}

\subsection*{Exercice 6 - Énergie électromagnétique dans un diélectrique \niveau{\moyen}}

\begin{boiteexercice}
Un condensateur plan de surface $S$ et de distance $d$ est rempli d'un diélectrique de permittivité $\epsilon_r$. Il est chargé sous une tension $U$.
\begin{enumerate}
  \item Calculer l'énergie électrostatique stockée.
  \item Calculer la densité d'énergie électromagnétique dans le diélectrique.
  \item Exprimer la force exercée sur le diélectrique si on l'introduit progressivement entre les armatures, à charge constante.
\end{enumerate}
\end{boiteexercice}

\begin{boitecorrige}

\textbf{1.} La capacité est $C = \epsilon_0\epsilon_r S/d$ et l'énergie stockée :
\[
W = \frac{1}{2}CU^2 = \frac{\epsilon_0\epsilon_r S U^2}{2d}.
\]

\textbf{2.} Le champ entre les armatures est $E = U/d$. La densité d'énergie dans le diélectrique est :
\[
u_e = \frac{1}{2}\vec{D}\cdot\vec{E} = \frac{1}{2}\epsilon_0\epsilon_r E^2 = \frac{\epsilon_0\epsilon_r U^2}{2d^2}.
\]
On vérifie bien que $W = u_e\cdot Sd$.

\textbf{3.} À charge $Q$ constante, l'énergie est $W = Q^2/(2C)$. En notant $x$ la fraction du condensateur remplie par le diélectrique (longueur $x$ sur la longueur totale $L$) :
\[
C(x) = \frac{\epsilon_0 S}{d}\!\left[\epsilon_r\,\frac{x}{L} + 1\cdot\frac{L-x}{L}\right] = C_0\!\left[1 + (\epsilon_r-1)\frac{x}{L}\right].
\]
La force est :
\[
F = -\frac{\partial W}{\partial x}\Bigg|_Q = \frac{Q^2}{2C^2}\frac{\partial C}{\partial x} = \frac{Q^2(\epsilon_r-1)}{2C_0 L}\cdot\frac{C_0^2}{C^2} \cdot \frac{\d C}{\d x}.
\]
Le diélectrique est aspiré vers l'intérieur du condensateur ($\epsilon_r > 1$, donc $F > 0$ dans le sens d'insertion) car cela diminue l'énergie à charge constante.

\end{boitecorrige}

\subsection*{Exercice 7 - Sphère diélectrique dans un champ uniforme \niveau{\difficile}}

\begin{boiteexercice}
Une sphère de rayon $R$, constituée d'un diélectrique LHI de permittivité relative $\epsilon_r$, est plongée dans un champ électrostatique uniforme $\vec{E}_0 = E_0\,\hat{e}_z$ dans le vide.
\begin{enumerate}
  \item Déterminer le potentiel $V_1$ à l'intérieur de la sphère et $V_2$ à l'extérieur, en cherchant des solutions à symétrie axiale de la forme $V(r,\theta) = \sum_l (A_l r^l + B_l r^{-(l+1)})P_l(\cos\theta)$.
  \item Exprimer le champ électrique $\vec{E}_\text{int}$ à l'intérieur de la sphère.
  \item Calculer la polarisation $\vec{P}$ et la densité surfacique de charges de polarisation $\sigma_p$.
\end{enumerate}
\end{boiteexercice}

\begin{boitecorrige}

\textbf{1. Potentiel à l'intérieur et à l'extérieur.}

À l'extérieur, le potentiel doit tendre vers $-E_0 r\cos\theta$ (champ uniforme) lorsque $r\to\infty$. À l'intérieur, il doit être fini en $r=0$. On conserve uniquement le terme $l=1$ :
\[
V_1(r,\theta) = -A\,r\cos\theta \quad (r<R),\qquad V_2(r,\theta) = -E_0\,r\cos\theta + \frac{B\cos\theta}{r^2} \quad (r>R).
\]
Les conditions aux limites à $r=R$ sont la continuité du potentiel et celle de la composante normale de $\vec{D} = \epsilon_0\epsilon_r\vec{E}$ :
\begin{align*}
V_1(R,\theta) &= V_2(R,\theta) \quad\Rightarrow\quad -AR = -E_0 R + B/R^2,\\
\epsilon_r\,\frac{\partial V_1}{\partial r}\Big|_R &= \frac{\partial V_2}{\partial r}\Big|_R \quad\Rightarrow\quad -\epsilon_r A = -E_0 - 2B/R^3.
\end{align*}
La résolution donne :
\[
A = \frac{3E_0}{\epsilon_r+2}, \qquad B = R^3 E_0\,\frac{\epsilon_r-1}{\epsilon_r+2}.
\]

\textbf{2. Champ intérieur.}

Le potentiel $V_1 = -A\,z$ correspond à un champ uniforme :
\[
\vec{E}_\text{int} = A\,\hat{e}_z = \frac{3}{\epsilon_r+2}\,\vec{E}_0.
\]
Le champ à l'intérieur est donc uniforme et réduit par le facteur $3/(\epsilon_r+2)$ : c'est le \emph{champ dépolarisant}.

\textbf{3. Polarisation et charges de polarisation.}

La polarisation est uniforme :
\[
\vec{P} = \epsilon_0(\epsilon_r-1)\vec{E}_\text{int} = 3\epsilon_0\,\frac{\epsilon_r-1}{\epsilon_r+2}\,\vec{E}_0.
\]
La densité surfacique de charges de polarisation est $\sigma_p = \vec{P}\cdot\hat{n} = P\cos\theta$, soit une distribution dipolaire surfacique. Le moment dipolaire total vaut $\vec{p} = \vec{P}\cdot(4\pi R^3/3)$, qui co\"incide avec le dipôle équivalent au champ extérieur ($B = p/(4\pi\epsilon_0)$).

\reponse{$\vec{E}_\text{int} = \dfrac{3}{\epsilon_r+2}\,\vec{E}_0$}
\end{boitecorrige}

\subsection*{Exercice 8 - Solénoïde torique à noyau ferromagnétique \niveau{\moyen}}

\begin{boiteexercice}
Un solénoïde torique de rayon moyen $R$ et de section circulaire de rayon $a\ll R$ comporte $N$ spires. Il est bobiné autour d'un noyau ferromagnétique de perméabilité $\mu_r = 1000$.
\begin{enumerate}
  \item Déterminer le champ $\vec{H}$, puis $\vec{B}$ à l'intérieur du tore en fonction du courant $I$.
  \item Calculer l'inductance $L$ de la bobine.
  \item On ouvre une mince fente d'épaisseur $e\ll 2\pi R$ dans le tore. En supposant que $\vec{B}$ reste tangentiel et continu, calculer le nouveau champ $H_\text{fer}$ dans le fer et $H_\text{air}$ dans la fente. Conclure.
\end{enumerate}
\end{boiteexercice}

\begin{boitecorrige}

\textbf{1. Champ dans le tore.}

Par symétrie, on applique le théorème d'Ampère sur un cercle de rayon $R$ :
\[
\oint \vec{H}\cdot\d\vec{\ell} = H\cdot 2\pi R = NI \quad\Rightarrow\quad H = \frac{NI}{2\pi R}.
\]
Le champ $B$ vaut alors :
\[
B = \mu_0\mu_r H = \frac{\mu_0\mu_r N I}{2\pi R}.
\]

\textbf{2. Inductance.}

Le flux à travers une spire est $\Phi_\text{spire} = B\cdot\pi a^2$. Le flux total pour $N$ spires :
\[
\Phi = N\Phi_\text{spire} = \frac{\mu_0\mu_r N^2 \pi a^2 I}{2\pi R}.
\]
L'inductance est $L = \Phi/I$, soit :
\[
L = \frac{\mu_0\mu_r N^2 a^2}{2R}.
\]

\textbf{3. Fente d'air.}

La continuité de la composante normale de $\vec{B}$ donne $B_\text{fer} = B_\text{air} = B$. En revanche, $H$ est discontinu :
\[
H_\text{fer} = \frac{B}{\mu_0\mu_r}, \qquad H_\text{air} = \frac{B}{\mu_0}.
\]
Le théorème d'Ampère devient :
\[
NI = H_\text{fer}(2\pi R - e) + H_\text{air}\,e = \frac{B}{\mu_0}\!\left[\frac{2\pi R - e}{\mu_r} + e\right].
\]
Comme $e\ll 2\pi R$ mais $\mu_r\gg 1$, le terme $e$ peut dominer. Par exemple pour $\mu_r = 1000$ et $e/(2\pi R) = 10^{-3}$, le rapport des champs vaut $H_\text{air}/H_\text{fer} = \mu_r$, et la fente d'air \emph{« désaimante »} fortement le noyau : le champ $B$ chute d'un facteur $\approx 2$. C'est pourquoi les circuits magnétiques réels doivent avoir des entrefers les plus courts possibles.

\reponse{$L = \dfrac{\mu_0\mu_r N^2 a^2}{2R}$}
\end{boitecorrige}

\subsection*{Exercice 9 - Capacité d'un câble coaxial à deux diélectriques \niveau{\moyen}}

\begin{boiteexercice}
Un câble coaxial est constitué d'un conducteur interne de rayon $a$ et d'une gaine externe de rayon interne $b$. L'espace entre les conducteurs est partiellement rempli de deux diélectriques concentriques : un diélectrique de permittivité $\epsilon_1$ entre $a$ et $c$ (avec $a<c<b$), et un diélectrique $\epsilon_2$ entre $c$ et $b$. La longueur du câble est $\ell$.
\begin{enumerate}
  \item Déterminer $\vec{D}$, $\vec{E}_1$ et $\vec{E}_2$ pour une charge $\lambda$ par unité de longueur.
  \item Calculer la ddp $U$ entre les conducteurs.
  \item En déduire la capacité par unité de longueur $C/\ell$.
\end{enumerate}
\end{boiteexercice}

\begin{boitecorrige}

\textbf{1. Champs.}

Par symétrie cylindrique, $\vec{D} = D(r)\hat{e}_r$. Le théorème de Gauss sur un cylindre de rayon $r$ ($a<r<b$) et de longueur $\ell$ :
\[
\oint \vec{D}\cdot\d\vec{S} = D\cdot 2\pi r \ell = \lambda \ell \quad\Rightarrow\quad \vec{D} = \frac{\lambda}{2\pi r}\hat{e}_r.
\]
Les champs électriques dans chaque diélectrique :
\[
\vec{E}_1 = \frac{\lambda}{2\pi\epsilon_1 r}\hat{e}_r \;(a<r<c), \qquad \vec{E}_2 = \frac{\lambda}{2\pi\epsilon_2 r}\hat{e}_r \;(c<r<b).
\]

\textbf{2. Différence de potentiel.}

\[
U = \int_a^b E(r)\,\d r = \int_a^c \frac{\lambda\,\d r}{2\pi\epsilon_1 r} + \int_c^b \frac{\lambda\,\d r}{2\pi\epsilon_2 r} = \frac{\lambda}{2\pi}\!\left[\frac{\ln(c/a)}{\epsilon_1} + \frac{\ln(b/c)}{\epsilon_2}\right].
\]

\textbf{3. Capacité linéique.}

\[
\frac{C}{\ell} = \frac{\lambda}{U} = \frac{2\pi}{\dfrac{\ln(c/a)}{\epsilon_1} + \dfrac{\ln(b/c)}{\epsilon_2}}.
\]
Les deux diélectriques en série s'additionnent comme des capacités en série. Pour $\epsilon_1 = \epsilon_2 = \epsilon$, on retrouve $C/\ell = 2\pi\epsilon/\ln(b/a)$.

\reponse{$C/\ell = \dfrac{2\pi}{\dfrac{1}{\epsilon_1}\ln\dfrac{c}{a} + \dfrac{1}{\epsilon_2}\ln\dfrac{b}{c}}$}
\end{boitecorrige}

\subsection*{Exercice 10 - Pression de radiation sur un miroir parfait \niveau{\moyen}}

\begin{boiteexercice}
Une onde électromagnétique plane progressive monochromatique d'intensité $I$ (en $\text{W\,m}^{-2}$) arrive normalement sur un miroir parfaitement conducteur (réflectivité $R=1$).
\begin{enumerate}
  \item Rappeler la relation entre $I$ et le module du vecteur de Poynting moyen $\langle\Pi\rangle$.
  \item Calculer la pression de radiation $P_\text{rad}$ exercée sur le miroir par transfert de quantité de mouvement.
  \item Application : un laser de puissance $1\;\text{kW}$ focalisé sur $1\;\text{mm}^2$. Calculer $P_\text{rad}$.
\end{enumerate}
\end{boiteexercice}

\begin{boitecorrige}

\textbf{1. Intensité.}

Pour une onde plane dans le vide, $\vec{E} = E_0\cos(kz-\omega t)\hat{e}_x$, $\vec{B} = E_0/c\cos(kz-\omega t)\hat{e}_y$. Le vecteur de Poynting vaut :
\[
\vec{\Pi} = \frac{\vec{E}\times\vec{B}}{\mu_0} = \frac{E_0^2}{\mu_0 c}\cos^2(kz-\omega t)\,\hat{e}_z.
\]
En moyenne temporelle $\langle\cos^2\rangle = 1/2$, donc :
\[
I = \langle\Pi\rangle = \frac{E_0^2}{2\mu_0 c}.
\]

\textbf{2. Pression de radiation.}

Le miroir parfait réfléchit l'onde : la quantité de mouvement par photon passe de $+p$ à $-p$, soit une variation $\Delta p = 2p$. Pour $n$ photons par unité de volume et de temps, le flux de quantité de mouvement réfléchi est $2I/c$. La pression de radiation est donc :
\[
P_\text{rad} = \frac{2I}{c}.
\]

\textbf{3. Application.}

Pour $P_\text{laser} = 1\;\text{kW}$ sur $S = 1\;\text{mm}^2 = 10^{-6}\;\text{m}^2$ :
\[
I = \frac{10^3}{10^{-6}} = 10^9\;\text{W\,m}^{-2}, \qquad P_\text{rad} = \frac{2\times 10^9}{3\times 10^8} \approx 6{,}67\;\text{Pa}.
\]
Soit une force $F = P_\text{rad}\cdot S \approx 6{,}67\;\mu\text{N}$.

\reponse{$P_\text{rad} = \dfrac{2I}{c}$}
\end{boitecorrige}

\subsection*{Exercice 11 - Théorème de Clausius-Mossotti pour les gaz \niveau{\moyen}}

\begin{boiteexercice}
On considère un gaz de molécules polaires de moment dipolaire permanent $p_0$, à la température $T$ et de densité volumique $n$.
\begin{enumerate}
  \item Établir la polarisation moyenne de Langevin $\langle p_z\rangle = p_0\mathcal{L}(\beta)$ avec $\beta = p_0 E/(k_B T)$.
  \item Dans la limite $\beta\ll 1$, montrer que la permittivité relative suit la loi de Debye : $\epsilon_r - 1 \approx \dfrac{n p_0^2}{3\epsilon_0 k_B T}$.
  \item Application numérique pour la vapeur d'eau à $T=300\;\text{K}$ ($p_0 = 6{,}2\times 10^{-30}\;\text{C\,m}$, $n$ correspondant à $1\;\text{atm}$).
\end{enumerate}
\end{boiteexercice}

\begin{boitecorrige}

\textbf{1. Polarisation de Langevin.}

L'énergie d'un dipôle dans le champ $E$ est $U = -p_0 E\cos\theta = -k_B T\beta\cos\theta$. La probabilité d'orientation est donnée par Boltzmann : $\d P \propto \e^{-\beta\cos\theta}\sin\theta\,\d\theta\,\d\phi/(4\pi)$. La moyenne du cosinus :
\[
\langle\cos\theta\rangle = \frac{\displaystyle\int_0^\pi \cos\theta\,\e^{-\beta\cos\theta}\sin\theta\,\d\theta}{\displaystyle\int_0^\pi \e^{-\beta\cos\theta}\sin\theta\,\d\theta} = \coth\beta - \frac{1}{\beta} = \mathcal{L}(\beta),
\]
où $\mathcal{L}$ est la fonction de Langevin. La polarisation moyenne est donc :
\[
\langle p_z\rangle = p_0\,\mathcal{L}(\beta).
\]

\textbf{2. Limite $\beta\ll 1$ (champ faible ou haute température).}

Le développement limité $\mathcal{L}(\beta)\approx \beta/3$ à l'ordre 1 donne :
\[
\langle p_z\rangle \approx \frac{p_0\,\beta}{3} = \frac{p_0^2 E}{3k_B T}.
\]
La polarisation macroscopique est $P = n\langle p_z\rangle$, soit :
\[
P = \frac{n p_0^2}{3 k_B T}\,E \quad\Rightarrow\quad \chi_e = \frac{P}{\epsilon_0 E} = \frac{n p_0^2}{3\epsilon_0 k_B T}.
\]
Comme $\epsilon_r = 1 + \chi_e$ : \quad $\epsilon_r - 1 \approx \dfrac{n p_0^2}{3\epsilon_0 k_B T}$, ce qui est la loi de Debye.

\textbf{3. Application numérique.}

À $T=300\;\text{K}$ et $1\;\text{atm}$, $n = P/(k_B T) \approx 10^5/(1{,}38\times 10^{-23}\times 300) \approx 2{,}5\times 10^{25}\;\text{m}^{-3}$.
\[
\epsilon_r - 1 = \frac{2{,}5\times 10^{25}\times(6{,}2\times 10^{-30})^2}{3\times 8{,}85\times 10^{-12}\times 1{,}38\times 10^{-23}\times 300} \approx 7\times 10^{-4}.
\]
Ce résultat est cohérent avec la vapeur d'eau à l'état gazeux dilué (l'eau liquide a $\epsilon_r\approx 80$ à cause des interactions dipolaires fortes).

\reponse{$\epsilon_r - 1 \approx \dfrac{n p_0^2}{3\epsilon_0 k_B T}$}
\end{boitecorrige}

\subsection*{Exercice 12 - Onde évanescente et réflexion totale \niveau{\difficile}}

\begin{boiteexercice}
Une onde électromagnétique polarisée TM se propage dans un milieu d'indice $n_1$ et atteint l'interface avec un milieu d'indice $n_2 < n_1$ sous une incidence $\theta_i > \theta_c = \arcsin(n_2/n_1)$ (réflexion totale).
\begin{enumerate}
  \item Montrer que le vecteur d'onde transmis $\vec{k}_2$ est complexe et déterminer ses composantes.
  \item Définir l'onde évanescente et écrire le champ $\vec{E}_t$ dans le milieu 2.
  \item Calculer le vecteur de Poynting moyen $\langle\vec{\Pi}\rangle$ dans le milieu 2 et conclure.
\end{enumerate}
\end{boiteexercice}

\begin{boitecorrige}

\textbf{1. Vecteur d'onde transmis.}

La conservation de la composante tangentielle $k_x = k_1\sin\theta_i$ et la relation de dispersion $|\vec{k}_2|^2 = n_2^2 \omega^2/c^2$ donnent :
\[
k_{2x} = k_1\sin\theta_i = \frac{n_1\omega}{c}\sin\theta_i, \qquad k_{2z}^2 = k_2^2 - k_{2x}^2 = \frac{\omega^2}{c^2}(n_2^2 - n_1^2\sin^2\theta_i).
\]
Comme $\theta_i > \theta_c$, on a $n_1\sin\theta_i > n_2$, donc $k_{2z}^2 < 0$. On pose :
\[
k_{2z} = \ii\kappa, \quad \kappa = \frac{\omega}{c}\sqrt{n_1^2\sin^2\theta_i - n_2^2} > 0.
\]

\textbf{2. Onde évanescente.}

Le champ transmis (polarisation TM) s'écrit :
\[
\vec{E}_t(\vec{r},t) = \vec{E}_{t0}\,\e^{\ii(k_{2x}x-\omega t)}\,\e^{-\kappa z}.
\]
L'onde « glisse » le long de l'interface (propagation selon $x$) tout en s'atténuant exponentiellement selon $z$ sur une profondeur caractéristique :
\[
\delta_\text{év} = \frac{1}{\kappa} = \frac{\lambda_0}{2\pi\sqrt{n_1^2\sin^2\theta_i - n_2^2}}.
\]

\textbf{3. Vecteur de Poynting moyen.}

Le champ magnétique $\vec{B}_t = \vec{k}_2\times\vec{E}_t/\omega$ a une composante complexe. Le vecteur de Poynting moyen s'écrit :
\[
\langle\vec{\Pi}\rangle = \frac{1}{2\mu_0}\text{Re}\!\left(\vec{E}_t\times\vec{B}_t^*\right).
\]
On montre que $\langle\Pi_z\rangle = 0$ : il n'y a pas de flux d'énergie en moyenne vers le milieu 2. En revanche, $\langle\Pi_x\rangle \neq 0$ : l'énergie se propage parallèlement à l'interface, à l'intérieur du milieu 2 sur la profondeur $\delta_\text{év}$. C'est l'\emph{onde évanescente}, utilisée notamment dans les microscopes de fluorescence totale interne (TIRF).

\reponse{$\delta_\text{év} = \dfrac{\lambda_0}{2\pi\sqrt{n_1^2\sin^2\theta_i - n_2^2}}$}
\end{boitecorrige}

\subsection*{Exercice 13 - Bilan énergétique de l'onde de plasma \niveau{\difficile}}

\begin{boiteexercice}
Dans un plasma froid, dilué et non collisionnel, les électrons (charge $-e$, masse $m_e$, densité $n_e$) sont mobiles devant des ions fixes. La permittivité diélectrique effective vaut :
\[
\epsilon(\omega) = \epsilon_0\!\left(1 - \frac{\omega_p^2}{\omega^2}\right), \qquad \omega_p^2 = \frac{n_e e^2}{\epsilon_0 m_e}.
\]
\begin{enumerate}
  \item Discuter la propagation d'une onde plane selon $\omega \lessgtr \omega_p$. Que vaut la vitesse de phase $v_\varphi$ et de groupe $v_g$ pour $\omega > \omega_p$ ?
  \item Pour $\omega < \omega_p$, montrer que l'onde est évanescente et calculer la profondeur de pénétration.
  \item Application : ionosphère terrestre ($n_e\approx 10^{12}\;\text{m}^{-3}$). Calculer $\omega_p$ et $f_p = \omega_p/(2\pi)$.
\end{enumerate}
\end{boiteexercice}

\begin{boitecorrige}

\textbf{1. Régime $\omega > \omega_p$ : propagation.}

La relation de dispersion est $k^2 c^2 = \omega^2 - \omega_p^2$, soit :
\[
k(\omega) = \frac{1}{c}\sqrt{\omega^2 - \omega_p^2}.
\]
L'indice de réfraction effectif vaut $n = \sqrt{1-\omega_p^2/\omega^2} < 1$. Les vitesses :
\[
v_\varphi = \frac{\omega}{k} = \frac{c}{\sqrt{1-\omega_p^2/\omega^2}} > c, \qquad v_g = \frac{\d\omega}{\d k} = c\sqrt{1-\frac{\omega_p^2}{\omega^2}} < c.
\]
On vérifie $v_\varphi v_g = c^2$ : aucun conflit avec la relativité (l'information voyage à $v_g < c$).

\textbf{2. Régime $\omega < \omega_p$ : onde évanescente.}

$k^2 < 0$, donc $k = \ii\kappa$ avec $\kappa = \sqrt{\omega_p^2 - \omega^2}/c$. L'onde est évanescente :
\[
\vec{E}(z,t) = \vec{E}_0\,\e^{-\kappa z}\,\e^{-\ii\omega t}.
\]
Profondeur de pénétration : $\delta = 1/\kappa = c/\sqrt{\omega_p^2-\omega^2}$. Pour $\omega\to 0$, $\delta\to c/\omega_p$ : c'est la \emph{longueur de pénétration inertielle}. Pour $\omega\to\omega_p^-$, $\delta\to\infty$ : l'onde pénètre très profondément juste sous la coupure.

\textbf{3. Application ionosphère.}

\[
\omega_p = \sqrt{\frac{n_e e^2}{\epsilon_0 m_e}} = \sqrt{\frac{10^{12}\times(1{,}6\times 10^{-19})^2}{8{,}85\times 10^{-12}\times 9{,}1\times 10^{-31}}} \approx 5{,}6\times 10^7\;\text{rad\,s}^{-1}.
\]
Soit $f_p \approx 9\;\text{MHz}$. Les ondes radio de fréquence inférieure à $f_p$ sont réfléchies par l'ionosphère (utilisées pour les communications longues distances par bonds), tandis que celles au-dessus de $f_p$ traversent (communications spatiales).

\reponse{$f_p = \dfrac{1}{2\pi}\sqrt{\dfrac{n_e e^2}{\epsilon_0 m_e}} \approx 9\;\text{MHz pour l'ionosphère}$}
\end{boitecorrige}

\ifdefined\inMainDoc\else\end{document}\fi
